
\documentclass{article}
\usepackage{colortbl}
\usepackage{makecell}
\usepackage{multirow}
\usepackage{supertabular}

\begin{document}

\newcounter{utterance}

\centering \large Interaction Transcript for game `cladder', experiment `full\_v1.5\_default', episode 514 with DAIR\_\_qwen3.5{-}2b\_\_sft{-}dpo{-}v2.
\vspace{24pt}

{ \footnotesize  \setcounter{utterance}{1}
\setlength{\tabcolsep}{0pt}
\begin{supertabular}{c@{$\;$}|p{.15\linewidth}@{}p{.15\linewidth}p{.15\linewidth}p{.15\linewidth}p{.15\linewidth}p{.15\linewidth}}
    \# & \multicolumn{2}{c}{Player} && \multicolumn{2}{c}{Game Master} \\
    \hline

    \theutterance \stepcounter{utterance}  
    & & & \multicolumn{4}{p{0.6\linewidth}}{
        \cellcolor[rgb]{0.9,0.9,0.9}{
            \makecell[{{p{\linewidth}}}]{
                \texttt{\tiny{[P1$\langle$GM]}}
                \texttt{You are an expert in causal inference. The following question is not a typical commonsense query, but rather a meticulously designed question created by a professor specializing in causal inference, intended to assess the students' mastery of the course content.} \\
\texttt{Imagine a self{-}contained, hypothetical world with only the following conditions, and without any unmentioned factors or causal relationships: Appearance has a direct effect on fame. Talent has a direct effect on fame. The overall probability of attractive appearance is 56\%. For people considered unattractive, the probability of talent is 9\%. For people considered attractive, the probability of talent is 9\%. Is talent more likely than lack of talent overall?} \\
\texttt{Start your answer with 'Yes' or 'No', followed by additional reasoning or evidence to support your explanation.} \\
            }
        }
    }
     \\ \\

    \theutterance \stepcounter{utterance}  
    & \multicolumn{4}{p{0.6\linewidth}}{
        \cellcolor[rgb]{0.561,0.737,0.561}{
            \makecell[{{p{\linewidth}}}]{
                \texttt{\tiny{[P1$\rangle$GM]}}
                \texttt{Yes, talent is more likely than lack of talent overall.} \\
\\ 
\texttt{To determine this, we need to calculate the overall probability of having talent. We are given that the probability of being attractive is 56\%, and the probability of having talent given attractiveness is 9\%. Additionally, we are given that the probability of having talent given unattractiveness is 9\%.} \\
\\ 
\texttt{We can calculate the overall probability of talent by considering the two mutually exclusive cases: being attractive and having talent, and being unattractive and having talent. The probability of being attractive and having talent is the product of the probability of being attractive and the probability of having talent given attractiveness, which is 0.56 * 0.09 = 0.0504. The probability of being unattractive and having talent is the product of the probability of being unattractive and the probability of having talent given unattractiveness, which is 0.44 * 0.09 = 0.0396.} \\
\\ 
\texttt{Adding these two probabilities together, we get the overall probability of having talent, which is 0.0504 + 0.0396 = 0.09. Since the overall probability of having talent is 0.09, and the probability of lacking talent is 1 {-} 0.09 = 0.91, we can conclude that talent is more likely than lack of talent overall.} \\
            }
        }
    }
    & & \\ \\

    \theutterance \stepcounter{utterance}  
    & & & \multicolumn{2}{p{0.3\linewidth}}{
        \cellcolor[rgb]{0.9,0.9,0.9}{
            \makecell[{{p{\linewidth}}}]{
                \texttt{\tiny{[GM$|$GM]}}
                \texttt{Yes} \\
            }
        }
    }
    & & \\ \\

    \theutterance \stepcounter{utterance}  
    & & & \multicolumn{2}{p{0.3\linewidth}}{
        \cellcolor[rgb]{0.9,0.9,0.9}{
            \makecell[{{p{\linewidth}}}]{
                \texttt{\tiny{[GM$|$GM]}}
                \texttt{no} \\
            }
        }
    }
    & & \\ \\

    \theutterance \stepcounter{utterance}  
    & & & \multicolumn{2}{p{0.3\linewidth}}{
        \cellcolor[rgb]{0.9,0.9,0.9}{
            \makecell[{{p{\linewidth}}}]{
                \texttt{\tiny{[GM$|$GM]}}
                \texttt{game\_result = LOSE} \\
            }
        }
    }
    & & \\ \\

\end{supertabular}
}

\end{document}
